\documentclass{beamer}
\usepackage[T1]{fontenc}
\usepackage[english]{babel}
\usepackage{lmodern}
\usepackage{amsmath,amssymb}
\usepackage{algorithm}
\usepackage{algorithmic}
\usepackage{graphicx}
\usepackage{subcaption}
\usepackage{caption}
\usepackage{tikz}
\usetikzlibrary{shapes.geometric,positioning,calc,fit}
\usetikzlibrary{overlay-beamer-styles}
\usepackage{booktabs,multicol,stackengine}
\usepackage[table,dvipsnames]{xcolor}
\usepackage{cancel}
\usepackage{moresize}
\usepackage{hyperref}

\usepackage{listings}
\definecolor{nus-orange}{RGB}{239,124,0}
\definecolor{nus-blue}{RGB}{0,61,124}
\definecolor{halfgray}{gray}{0.55}

\lstset{
  basicstyle=\ttfamily\scriptsize,
  keywordstyle=\bfseries\color{nus-blue},
  stringstyle=\color{ForestGreen},
  commentstyle=\color{halfgray}\itshape,
  numbers=left,
  numberstyle=\scriptsize\color{halfgray},
  rulesepcolor=\color{nus-orange},
  frame=shadowbox,
  columns=flexible,
  showstringspaces=false,
  keepspaces=true,
  tabsize=2,
  breaklines=true
}
% PL/pgSQL keywords on top of the SQL dialect
\lstdefinelanguage{plpgsql}[]{SQL}{
  morekeywords={RETURNS,TRIGGER,DECLARE,BEGIN,END,IF,THEN,ELSE,RETURN,
    RAISE,EXCEPTION,REPLACE,FUNCTION,EXECUTE,BEFORE,AFTER,EACH,ROW,
    CONSTRAINT,DEFERRABLE,INITIALLY,DEFERRED,LANGUAGE,RECORD,NEW,OLD,
    EXISTS,ISNULL,CASCADE,SCHEMA},
  morecomment=[l]{--},
  morecomment=[s]{/*}{*/},
}

\renewcommand{\figurename}{Figure}
\renewcommand{\algorithmname}{Algorithm}

\usepackage{SHU}
% VITAL STUFF
\renewcommand\epsilon\varepsilon
\renewcommand\phi\varphi
\renewcommand\leq\leqslant
\renewcommand\geq\geqslant
\newcommand\defeq{\stackrel{\mathrm{def}}{=}}

% For revision
\newcommand\rev[1]{{#1}}

% COMMENTS
\newcommand\mikael[1]{\textcolor{orange}{[\textbf{Mikaël:} #1]}}
\newcommand\pierre[1]{\textcolor{magenta}{[\textbf{Pierre:} #1]}}
\newcommand\pratik[1]{\textcolor{blue}{[\textbf{Pratik:} #1]}}

% SMOL SYMBOLS
\newcommand\calF{\mathcal{F}}
\newcommand\RR{\mathbb{R}}
\newcommand\QQ{\mathbb{Q}}
\newcommand\bD{\mathbf{D}}

% REDUCTIONS
\newcommand{\tr}{\leq_{\mathsf{P}}}
\newcommand{\equivt}{\equiv_{\mathsf{P}}}

% PROBLEMS
\newcommand\escore{\mathsf{EScore}}
\newcommand\score{\mathsf{Score}}
\newcommand\eshapley{\escore_{\cshapley}}
\newcommand\shapley{\score_{\cshapley}}
\newcommand\ebanzhaf{\escore_{\cbanzhaf}}
\newcommand\banzhaf{\score_{\cbanzhaf}}
\newcommand\ev{\mathsf{EV}}
\newcommand\evs{\mathsf{EV}_\star}
\newcommand\ennv{\mathsf{ENV}} % \env was already defined
\newcommand\envss{\mathsf{ENV}_{\star,\star}}
\newcommand\mc{\mathsf{MC}}
\newcommand\env{\mathsf{ENV}}

% NOTATIONS
\newcommand\vars{\mathsf{Vars}}
\newcommand\out{\mathsf{out}} % for output gate
\newcommand\cshapley{c_{\mathrm{Shapley}}}
\newcommand\cbanzhaf{c_{\mathrm{Banzhaf}}}
\newcommand\ar{\mathsf{ar}}
\newcommand\const{\mathsf{Const}}
\newcommand{\pqe}{\mathrm{PQE}}
\newcommand\prov{\mathrm{Prov}}

% For algorithm2e
\newcommand{\hrulealg}[0]{\vspace{1mm} \hrule \vspace{1mm}}

% For alignment of numbers
\newcommand{\0}{\phantom{0}}

% Example
\newcommand{\phiex}{\phi_{\mathrm{ex}}}


\title{IT5008: Tutorial 7 --- Triggers}
\author{\href{https://pratik2358.github.io/}{Pratik Karmakar}}
\institute{School of Computing,\\ National University of Singapore}
\date{AY26/27 S1}

\begin{document}

\begin{frame}
  \titlepage
  \IfFileExists{nus-logo.png}{
    \begin{figure}[htpb]\centering
      \includegraphics[keepaspectratio, scale=0.18]{nus-logo.png}
    \end{figure}
  }{}
\end{frame}

\section{Scenario}

\begin{frame}{The Book Exchange Schema (Recap)}
\small
Same NUN Book Exchange scenario as before; the discussion questions use:
\begin{itemize}
  \item \textbf{student}(name, \underline{email}, year, department, graduate) --- with \texttt{CHECK(graduate >= year)}
  \item \textbf{loan}(borrower, owner, book, copy, borrowed, returned) --- \texttt{returned IS NULL} means \emph{still out}
\end{itemize}
The preliminary questions use two small, separate schemas (employees; projects).
\medskip
\begin{block}{New this tutorial}
\textbf{Triggers}: PL/pgSQL functions that run automatically on \texttt{INSERT}/\texttt{UPDATE}/\texttt{DELETE} --- to enforce constraints that \texttt{CHECK}, keys and foreign keys cannot express, or to \emph{repair} a row instead of rejecting it.
\end{block}
\end{frame}

\begin{frame}{Trigger Cheat Sheet (Row-Level)}
\footnotesize
\begin{center}
\renewcommand{\arraystretch}{1.3}
\begin{tabular}{p{0.2\linewidth}p{0.34\linewidth}p{0.34\linewidth}}
\toprule
 & \textbf{\texttt{BEFORE}} & \textbf{\texttt{AFTER}} \\
\midrule
Runs & before the row is written & after the row is written \\
\texttt{RETURN NEW} & write the row (possibly \emph{modified}) & ignored \\
\texttt{RETURN NULL} & \textbf{silently skip} this row & ignored \\
\texttt{RAISE EXCEPTION} & abort the statement & abort the statement \\
Good for & reject or \textbf{fix up} one row & check the \textbf{resulting state} \\
\bottomrule
\end{tabular}
\end{center}
\begin{itemize}
  \item \texttt{NEW} / \texttt{OLD}: the row after / before the change.
  \item \texttt{CREATE CONSTRAINT TRIGGER \ldots\ DEFERRABLE INITIALLY DEFERRED}: an \texttt{AFTER} trigger that runs at \textbf{\texttt{COMMIT}}, so a transaction may pass through invalid intermediate states.
  \item A trigger function must reach a \texttt{RETURN} (or raise) --- even when the value is ignored.
\end{itemize}
\end{frame}

\section{Preliminary Questions}

\begin{frame}{Preliminary: Triggers}
\small
\textit{Not discussed in class --- attempt before the tutorial, since later questions build on these.}
\begin{enumerate}
  \item[\textbf{1.}] \textbf{Triggers.}
  \begin{enumerate}
    \item[(a)] \texttt{Employees(\underline{eid}, name)}, \texttt{Engineers(\underline{eid}, skill)}, \texttt{Managers(\underline{eid}, dept)}; both subclasses reference \texttt{Employees.eid}.\\
    No employee can be \textbf{both} an engineer and a manager.
    \item[(b)] \texttt{Engineers(\underline{eid}, skill)}, \texttt{Projects(\underline{pid}, max\_hours)}, \texttt{Work(\underline{eid, pid}, hour)}; \texttt{Work} references both.\\
    The total \texttt{hour} worked on a project cannot exceed its \texttt{max\_hours}.
  \end{enumerate}
\end{enumerate}
\medskip
Enforce each with triggers.
\end{frame}

\begin{frame}[fragile,shrink=6]{1(a). Engineers Must Not Be Managers}
\footnotesize
{\scriptsize \textbf{Schema:} \texttt{Employees(\underline{eid}, name)}, \texttt{Engineers(\underline{eid}, skill)}, \texttt{Managers(\underline{eid}, dept)}\\
\textbf{FKs:} \texttt{Engineers.eid} $\to$ \texttt{Employees.eid}, \ \texttt{Managers.eid} $\to$ \texttt{Employees.eid}}
\medskip

A \texttt{BEFORE} trigger that \textbf{skips} the row (\texttt{RETURN NULL}) if the employee is already a manager.
\begin{lstlisting}[language=plpgsql]
CREATE OR REPLACE FUNCTION not_manager()
RETURNS TRIGGER AS $$
DECLARE
  count NUMERIC;
BEGIN
  SELECT COUNT(*) INTO count
  FROM Managers
  WHERE NEW.eid = Managers.eid;
  IF count > 0 THEN
    RETURN NULL;
  ELSE
    RETURN NEW;
  END IF;
END;
$$ LANGUAGE plpgsql;

CREATE TRIGGER non_manager
BEFORE INSERT OR UPDATE ON Engineers
FOR EACH ROW EXECUTE FUNCTION not_manager();
\end{lstlisting}
\end{frame}

\begin{frame}[fragile,shrink=6]{1(a). \ldots and Managers Must Not Be Engineers}
\footnotesize
The reference solution leaves this as an exercise. Without it, \texttt{INSERT INTO Managers} for an existing engineer still \textbf{succeeds}.
\begin{lstlisting}[language=plpgsql]
CREATE OR REPLACE FUNCTION not_engineer()
RETURNS TRIGGER AS $$
BEGIN
  IF EXISTS (SELECT * FROM Engineers e WHERE e.eid = NEW.eid) THEN
    RETURN NULL;
  ELSE
    RETURN NEW;
  END IF;
END;
$$ LANGUAGE plpgsql;

CREATE TRIGGER non_engineer
BEFORE INSERT OR UPDATE ON Managers
FOR EACH ROW EXECUTE FUNCTION not_engineer();
\end{lstlisting}
\begin{itemize}
  \item A constraint across two tables needs a trigger on \textbf{each} table that can break it.
  \item \texttt{RETURN NULL} rejects \emph{silently} (\texttt{INSERT 0 0}); use \texttt{RAISE EXCEPTION} if the caller should know.
\end{itemize}
\end{frame}

\begin{frame}[fragile,shrink=4]{1(b). Project Hours Must Not Exceed \texttt{max\_hours}}
\footnotesize
{\scriptsize \textbf{Schema:} \texttt{Engineers(\underline{eid}, skill)}, \texttt{Projects(\underline{pid}, max\_hours)}, \texttt{Work(\underline{eid, pid}, hour)}\\
\textbf{FKs:} \texttt{Work.eid} $\to$ \texttt{Engineers.eid}, \ \texttt{Work.pid} $\to$ \texttt{Projects.pid}\\
\textbf{Rule:} for each project, \texttt{SUM(hour)} over its \texttt{Work} rows $\le$ \texttt{max\_hours}.}
\begin{lstlisting}[language=plpgsql]
CREATE OR REPLACE FUNCTION not_exceeded()
RETURNS TRIGGER AS $$
BEGIN
  IF EXISTS (
    SELECT 1
    FROM Projects p, Work w
    WHERE p.pid = w.pid
    GROUP BY p.pid
    HAVING SUM(w.hour) > p.max_hours
  ) THEN
    RAISE EXCEPTION 'Hours exceeded!';
  END IF;
  RETURN NULL;
END;
$$ LANGUAGE plpgsql;
\end{lstlisting}
\textbf{Look for violations}: any project whose total hours exceed its limit.
\end{frame}

\begin{frame}[fragile]{1(b). The Triggers}
\small
\begin{lstlisting}[language=plpgsql]
CREATE CONSTRAINT TRIGGER project_limit
AFTER INSERT OR UPDATE ON Projects
DEFERRABLE INITIALLY DEFERRED
FOR EACH ROW EXECUTE FUNCTION not_exceeded();

CREATE CONSTRAINT TRIGGER work_limit
AFTER INSERT OR UPDATE ON Work
DEFERRABLE INITIALLY DEFERRED
FOR EACH ROW EXECUTE FUNCTION not_exceeded();
\end{lstlisting}
\begin{itemize}
  \item \texttt{CONSTRAINT TRIGGER} + \texttt{DEFERRABLE INITIALLY DEFERRED}: the check runs at \texttt{COMMIT}.
  \item One trigger per table that can break the rule (next slide: why).
\end{itemize}
\end{frame}

\begin{frame}[fragile]{1(b). Why Two Deferred Triggers?}
\small
\begin{itemize}
  \item \textbf{Both tables can break it:} adding/updating \texttt{Work} raises the total; updating \texttt{Projects} can lower \texttt{max\_hours} below it.
  \item \textbf{Deferred to \texttt{COMMIT}:} the check looks at the \emph{whole} database, so we let a transaction finish its changes before checking.
\end{itemize}
\begin{lstlisting}[language=plpgsql]
BEGIN;
INSERT INTO Projects VALUES (1, 10);
INSERT INTO Work VALUES (1, 1, 5);
COMMIT;                -- OK: 5 <= 10

BEGIN;
INSERT INTO Projects VALUES (2, 10);
INSERT INTO Work VALUES (1, 2, 50);
COMMIT;                -- ERROR: Hours exceeded!

UPDATE Projects SET max_hours = 4 WHERE pid = 1;
                       -- ERROR: Hours exceeded!
\end{lstlisting}
\end{frame}

\section{Discussion Questions}

\begin{frame}{Discussion: Triggers on the Book Exchange}
\small
\textit{To be discussed in class --- attempt before the tutorial; you may be called on to present.}
\begin{enumerate}
  \item[\textbf{2.}] \textbf{Triggers.}
  \begin{enumerate}
    \item[(a)] Adeline Wong (\texttt{awong007@msn.com}) has already borrowed 6 books and returned none.\\
    New constraint: a student may borrow \textbf{at most 6 books at a time} --- with 6 unreturned books, they cannot borrow another.
    \item[(b)] \texttt{student} has \texttt{CHECK(graduate >= year)}. We also want the graduation date to \textbf{not be in the future} --- but instead of rejecting a bad value (as \texttt{CHECK(graduate <= CURRENT\_DATE)} would), \textbf{change it} to the current date. Invoke on insert and update.
  \end{enumerate}
\end{enumerate}
\end{frame}

\begin{frame}[fragile]{2(a). At Most 6 Books at a Time --- Approach}
\small
Easiest with an \texttt{AFTER} trigger: let the row in, then \textbf{look for violating rows; if any, raise}.
\begin{lstlisting}[language=plpgsql,basicstyle=\ttfamily\footnotesize]
IF EXISTS ( /* SQL query checking violating rows */ ) THEN
  RAISE EXCEPTION '..';
END IF;
\end{lstlisting}
\begin{itemize}
  \item \textbf{Violating row} here: a borrower with \textbf{more than 6} loans where \texttt{returned IS NULL}.
  \item Raising inside an \texttt{AFTER} trigger aborts the statement, so the offending loan is never kept.
\end{itemize}
\end{frame}

\begin{frame}[fragile]{2(a). The Trigger Function}
\footnotesize
\begin{lstlisting}[language=plpgsql]
CREATE OR REPLACE FUNCTION loan_limit()
RETURNS TRIGGER AS $$
DECLARE
  violating_student RECORD;
BEGIN
  -- Is there any student with more than 6 active loans?
  SELECT l.borrower INTO violating_student
  FROM loan l
  WHERE l.returned ISNULL
  GROUP BY l.borrower
  HAVING COUNT(*) > 6;

  IF violating_student IS NOT NULL THEN
    RAISE EXCEPTION '% has borrowed more than 6 books',
      violating_student;
  ELSE
    RETURN NEW;
  END IF;
END;
$$ LANGUAGE plpgsql;
\end{lstlisting}
\end{frame}

\begin{frame}[fragile]{2(a). The Trigger}
\small
\begin{lstlisting}[language=plpgsql]
CREATE TRIGGER loan_limit_trigger
AFTER INSERT OR UPDATE ON loan
FOR EACH ROW EXECUTE FUNCTION loan_limit();
\end{lstlisting}
\begin{itemize}
  \item \textbf{\texttt{> 6}, not \texttt{>= 6}:} the check runs \emph{after} the new row is in. Adeline's existing 6 are fine; a 7th makes the count 7 and is rejected.
  \item \textbf{\texttt{UPDATE} too:} setting \texttt{returned} back to \texttt{NULL} re-opens a loan.
  \item Once she returns a book, she can borrow again.
\end{itemize}
\medskip
To remove it:
\begin{lstlisting}[language=plpgsql]
DROP TRIGGER loan_limit_trigger ON loan;
DROP FUNCTION loan_limit();
-- or wipe everything:
DROP SCHEMA public CASCADE; CREATE SCHEMA public;
\end{lstlisting}
\end{frame}

\begin{frame}[fragile]{2(a). Alternative: a \texttt{BEFORE} Trigger}
\footnotesize
Check only the \textbf{new borrower}, \emph{before} the row is written --- so the limit is \texttt{>= 6}.
\begin{lstlisting}[language=plpgsql]
CREATE OR REPLACE FUNCTION loan_limit_before()
RETURNS TRIGGER AS $$
BEGIN
  IF NEW.returned IS NULL AND (
       SELECT COUNT(*) FROM loan l
       WHERE l.borrower = NEW.borrower
         AND l.returned IS NULL) >= 6
  THEN
    RAISE EXCEPTION '% already has 6 books', NEW.borrower;
  END IF;
  RETURN NEW;
END;
$$ LANGUAGE plpgsql;

CREATE TRIGGER loan_limit_before_trigger
BEFORE INSERT ON loan
FOR EACH ROW EXECUTE FUNCTION loan_limit_before();
\end{lstlisting}
\textbf{Off-by-one:} \texttt{BEFORE} sees the table \emph{without} the new row, \texttt{AFTER} sees it \emph{with} it.
\end{frame}

\begin{frame}[fragile,shrink=3]{2(b). No Future Graduation --- Fix, Don't Reject}
\footnotesize
An \texttt{AFTER} trigger can only reject. To \textbf{change} the row, we need \texttt{BEFORE} and a modified \texttt{NEW}.
\begin{lstlisting}[language=plpgsql]
CREATE OR REPLACE FUNCTION graduation_date()
RETURNS TRIGGER AS $$
BEGIN
  IF NEW.graduate > CURRENT_DATE THEN
    NEW.graduate := CURRENT_DATE;
  END IF;
  RETURN NEW;
END;
$$ LANGUAGE plpgsql;

CREATE TRIGGER no_future_graduation
BEFORE INSERT OR UPDATE ON student
FOR EACH ROW EXECUTE FUNCTION graduation_date();
\end{lstlisting}
\begin{itemize}
  \item \texttt{UPDATE student SET graduate = '2030-01-01' \ldots} stores today's date instead.
  \item \textbf{Still can fail} when \texttt{CURRENT\_DATE < year}: the clamped date violates \texttt{CHECK(graduate >= year)}.
\end{itemize}
\end{frame}

\section{Guidelines}

\begin{frame}{Guidelines \& Common Pitfalls}
\footnotesize
\begin{itemize}\itemsep2pt
  \item \textbf{\texttt{BEFORE} to fix or skip a row, \texttt{AFTER} to check the result.} Only \texttt{BEFORE} can change \texttt{NEW} (2(b)).
  \item \textbf{Cover every table that can break the rule} --- both subclasses in 1(a), both \texttt{Work} and \texttt{Projects} in 1(b).
  \item \textbf{Query for violations, not for valid rows.} \texttt{IF EXISTS(\ldots) THEN RAISE} only works if the query finds the \emph{bad} rows (1(b)).
  \item \textbf{Always reach a \texttt{RETURN}}, even in \texttt{AFTER} triggers (1(b)).
  \item \textbf{Mind the off-by-one:} \texttt{AFTER} counts the new row, \texttt{BEFORE} does not (2(a)).
  \item \textbf{Defer whole-database checks} with \texttt{CONSTRAINT TRIGGER \ldots\ DEFERRABLE INITIALLY DEFERRED} when a transaction needs several steps to become valid (1(b)).
  \item \textbf{Silent vs.\ loud:} \texttt{RETURN NULL} in \texttt{BEFORE} drops the row without an error; \texttt{RAISE EXCEPTION} aborts it.
\end{itemize}
\end{frame}

\begin{frame}
\begin{center}
Questions?\\
Drop a mail at: pratik.karmakar@u.nus.edu
\end{center}
\end{frame}

\end{document}
